Define nonuniformly learnable classes and state related theorems 

A hypothesis class \(\mathcal{H}\) is nonuniformly learnable if there exist a learning algorithm \( A \), and a function
\(m_{\mathcal{H}}^{\mathrm{NUL}} : (0,1)^2 \times \mathcal{H} \to \mathbb{N}\)
such that, for every \(\epsilon, \delta \in (0,1)\) and for every \(h \in \mathcal{H}\), 
if \(m \geq m_{\mathcal{H}}^{\mathrm{NUL}}(\epsilon, \delta, h)\) then for every distribution \(\mathcal{D}\), 
with probability at least \(1 - \delta\) over the choice of \(S \sim \mathcal{D}^m\), it holds that
\[
  L_{\mathcal{D}}(A(S)) \leq L_{\mathcal{D}}(h) + \epsilon.
\]
Note that agnostic learnability implies nonuniformly learnability.

Theorem 7.2 
A hypothesis class \( \mathcal{H} \) of binary classifiers is nonuniformly learnable 
if and only if it is a countable union of agnostic PAC learnable hypothesis classes.

Theorem 7.3
Let \( \mathcal{H} \) be a hypothesis class that can be written as a countable  union of hypothesis classes, 
\( \mathcal{H} = \cup_{n \in \mathbb{N}} \mathcal{H}_n\) , where each \( \mathcal{H}_n \) 
enjoys the uniform convergence property. Then, \( \mathcal{H} \) is nonuniformly learnable.

Define the associated function
\[
  \varepsilon_n(m, \delta) = \inf \{ \varepsilon \in (0, 1) : m_{\mathcal{H}_n}^{\text{UC}}(\varepsilon, \delta) \leq m \}  
\]
that is the lower bound for the gap between empirical and true risks achievable using \( m \) samples.
Consider a weight function \( w : \mathbb{N} \to [0, 1] \) with \( \sum_{n=1}^\infty w(n) \leq 1 \),
which introduces a preference (or a priori bias) over the hypothesis classes \( \{ \mathcal{H}_n \}_{n \in \mathbb{N}} \).

Theorem 7.4
Suppose each of the \( \mathcal{H}_n \) has the uniform convergence property.
Then, for every \( \delta \in (0, 1) \) and distribution \( \mathcal{D} \), 
with probability of at least \( 1 - \delta \) over the choice of \( S \sim \mathcal{D}^m \), 
the following bound holds (simultaneously) for every \( n \in \mathbb{N} \) and \( h \in \mathcal{H}_n \):
\[ 
| L_{\mathcal{D}}(h) - L_S(h) | \leq \varepsilon_n(m, w(n) \cdot \delta)
\]
In particular it holds with probability \( 1 - \delta \) and for all \( h \in \cup_{n \in \mathbb{N}} \mathcal{H}_n \)
\[
 L_{\mathcal{D}} (h) \leq L_S(h) + \min_{n,h \in \mathcal{H}_n} \varepsilon_n(m, w(n)\cdot \delta)
\]

This motivates the following algorithm, let \( n(h) = \min \{ n : h \in \mathcal{H}_n \} \)
then finding the solution to  \( \argmin_{h \in \mathcal{H}} L_S(h) + \varepsilon_{n(h)}(m, w(n(h))\cdot \delta) \) 
is called Structural Risk Minimization (SRM). 
Intuitively we trade off some empirical risk minimization with our preference for certain \( \mathcal{H}_n \) via \( w \).

Theorem
Choose \( w(n) = \frac{6}{n^2 \pi^2} \). For \( \mathcal{H} \) as above, \( \mathcal{H} \) is nonuniformly learnable using
the SRM algorithm with rate
\[ 
  m_{\mathcal{H}}^{\text{NUL}}(\varepsilon, \delta, h) \leq m_{\mathcal{H}_{n(h)}}^{\text{UC}}(\frac{\varepsilon}{2}, \frac{6\delta}{(\pi n(h))^2})
\]


Define the minimum description length rule and consistency

Let \( \mathcal{H} \) be a hypothesis class, and \( \Sigma \) some alphabet with a function \( d : \mathcal{H} \to \Sigma^\ast \)
(the codomain are all words in terms of this alphabet). Call \( |h| = |d(h)| \) the length of the description \( d(h) \) for \( h \).
A language is prefix-free, if and only if for no distinct \( h_1, h_2 \), \( d(h_1) \) are the first \( |h_1| \) symbols of \( d(h_2) \).

Theorem
Let \( \mathcal{H} \) be a hypothesis class and let \( d : H \to \{0, 1\}^\ast \) be a prefix-free 
description language for H. Then, for every sample size, \( m \), every confidence
parameter, \( \delta > 0 \), and every probability distribution, \( \mathcal{D} \), with probability greater
than \( 1 - \delta \) over the choice of \( S \tilde \mathcal{D}^m \) we have for all \( h \in \mathcal{H} \),
\[
  L_{\mathcal{D}}(h) \leq L_{S}(h) + \sqrt{\frac{|h| + \ln(\frac{2}{\delta})}{2m}}
\]

The minimum description length algorithm then solves the problem
\[
  \argmin_{h \in \mathcal{H}} L_{S}(h) + \sqrt{\frac{|h| + \ln(\frac{2}{\delta})}{2m}}
\]

Note how this implies for hypotheses \( h_1, h_2 \) with same empirical risk, the one with shorter  description has better true risk minimization.
This could be considered to be an instance of Occam's razor. 
The clue is here that a choice of \( d \) is similar to a choice of some preference function \( w \).

Consistency

Let \( Z \) be a domain set, let \( \mathcal{P} \) be a set of
probability distributions over \( Z \), and let \( \mathcal{H} \) be a hypothesis class. 
A learning rule \( A \) is consistent with respect to \( \mathcal{h} \) and \( \mathcal{P} \) 
if there exists a function \( m^{\text{CON}}_{\mathcal{H}} : (0, 1)^2 \times \mathcal{H} \times \mathcal{P} \to \mathbb{N} \) s.t.,
for every \( \varepsilon, \delta (0, 1) \) every \( h \in \mathcal{H} \) and every \( \mathcal{D} \in \mathcal{P} \), if
\( m \geq  m^{\text{CON}}_{\mathcal{H}}(\varepsilon, \delta, h, \mathcal{D}) \) then with probability of at least \( 1 - \delta \) over
the choice of \( S \sim \mathcal{D}^m \) it holds that
\[
  L_{\mathcal{D}}(A(S)) \leq L_{\mathcal{D}}(h) + \varepsilon
\]
If \( \mathcal{P} \) is the set of all distributions, we say that \( A \) is universally consistent with
respect to \( \mathcal{H} \).

Consistency also includes flexibility in the underlying choice of distribution and is the weakest of all the learnability conditions.
Unlike the notions of PAC learnability and nonuniform learnability, the definition of consistency does not yield a 
natural learning algorithm or a way to encode prior knowledge. 
In fact, in many cases there is no need for prior knowledge at all.

We can make any nonuniform learner for \( \mathcal{H} \) be consistent with respect to the class of all classifiers \( \mathcal{X} \to \mathcal{Y} \):
Upon receiving a training set, we will first run the nonuniform learner over the training set, and then we will obtain a bound
on the true risk of the learned predictor. If this bound is small enough we are
done. Otherwise, we revert to the Memorize algorithm (simply returning the most common label of the training set for the input). 
This simple modification makes the algorithm consistent with respect to all classifiers.
Thus it is not of interest to prefer any consistent learner but has no stronger learnability condition over one which is 
only nonuniformly or PAC learnable.

In general PAC learning is quite preferable because it gives a strict answer for the number of samples required.
